\paragraph{Reflection positivity of the Euclidean measure at \(N=2\): failure on the
invariant algebra and the even-subalgebra bound (conditional; bounded no-go).}
Fix the reflection \(\theta=\mathrm{diag}(-1,1,\dots,1)\), acting on the Euclidean
variables of \(\mathcal R_1\) by \(\tilde A_0\mapsto-\tilde A_0\),
\(\tilde A_i\mapsto\tilde A_i\), and consider the normalized reflected pairing
\(\langle f,f\rangle_\theta=\int f^\theta\bar f\,\rho\,\mathrm d\tilde A\big/
\int\rho\,\mathrm d\tilde A\) on polynomial observables, with
\(f^\theta(\tilde A)=f(\theta\tilde A)\) and \(\rho\) one of the Euclidean weights on
the right-hand side of the displayed limit identity: \(\rho=c_{\mathrm P}
(\det\tilde{\mathcal G})^4\mathrm e^{-S_{\mathrm E}}\) at \(N=2\) with the content
above (integrable by the paragraph on the bound along the flat directions), the
Pfaffian reading there following from \(\operatorname{Pf}\mathcal M=c_{\mathrm P}
(\det\mathcal G)^4\) with \(\mathcal G_{ab}=\eta^{\mu\nu}A^a_\mu A^b_\nu\) at
\(A=(-\mathrm i\tilde A_0,\tilde A_i)\), or \(\rho=\mathrm e^{-S_{\mathrm E}}\mathrm
e^{-\frac12\sum_{\mu,\nu}M_{\mu\nu}\operatorname{Tr}\tilde A_\mu\tilde A_\nu}\) for
any positive definite block-diagonal \(M\) with \(M_{0i}=0\) at any \(N\); constant
factors cancel in the ratio, and the denominator is nonzero for the Pfaffian weight
at \(N=2\) and strictly positive with a mass. Each commutator in \(S_{\mathrm E}\)
changes sign exactly when it carries one time index, so \(S_{\mathrm E}\), the
Lebesgue measure, every such mass factor and the Gram
\(\tilde{\mathcal G}_{ab}=\sum_\mu\tilde A^a_\mu\tilde A^b_\mu\) are
\(\theta\)-invariant, and with \(f_\pm=(f\pm f^\theta)/2\) one has
\(\langle f,f\rangle_\theta=(\int|f_+|^2w-\int|f_-|^2w)/\int w\) with \(w\ge0\) the
real weight. Hence every \(\theta\)-odd polynomial not vanishing almost everywhere
has strictly negative reflected value \(-\int|f|^2w/\int w\), and such polynomials
are gauge invariant: \(\operatorname{Tr}\tilde A_0\tilde A_i=
\frac12\sum_a\tilde A^a_0\tilde A^a_i\), and, with
\(A_\mu=\sum_aA^a_\mu\sigma^a/2\), \(\operatorname{Tr}([\tilde A_0,\tilde A_i]
\tilde A_j)=\frac{\mathrm i}2\epsilon^{abc}\tilde A^a_0\tilde A^b_i\tilde A^c_j\) for
\(i\ne j\); both are nonzero and \(\theta\)-odd and invariant by cyclicity of the
trace (with the Gram entries they are the generators of the invariant polynomial
algebra of adjoint tuples at \(N=2\), Weyl's first fundamental theorem for
\(\mathrm{SO}(3)\) with the \(\epsilon\)-contractions). Therefore the measures of
this weight class, with \(p=4\) content at \(N=2\) and with or without a mass, and
bosonically with a mass at every finite \(N\), do not satisfy the reflection
positivity of~\cite{OsterwalderSchrader1973,OsterwalderSchrader1975} on the
invariant polynomial algebra: the failure sits on manifestly gauge-invariant
bilinear and cubic singlets, not only on the translation labels of
Section~\ref{sec:status-D}. Any unital \(*\)-subalgebra closed under \(\theta\) on which \(\langle\cdot,\cdot\rangle_\theta\) is nonnegative contains no nonzero \(\theta\)-odd element, since such an element would have negative square, so every \(\theta\)-stable reflection-positive test algebra for these measures lies in the \(\theta\)-even invariant algebra, and every invariant mixing time and space under \(\theta\) is excluded from any reconstruction from these measures. On the
continuation side \(\operatorname{Tr}A_0A_i\) restricted to \(\mathcal R_1\) is
\(-\frac{\mathrm i}2\sum_a\tilde A^a_0\tilde A^a_i\) in the Pfaffian variables
\((-\mathrm i\tilde A_0,\tilde A_i)\) and \(\frac{\mathrm e^{-\pi\mathrm
i/4}}2\sum_a\tilde A^a_0\tilde A^a_i\) on the observable contour, still
\(\theta\)-odd up to a constant phase of unit modulus with strictly negative
reflected value: the negative sign passes to the Euclidean-side continuation of the
mixed two-point function. These statements are for the Euclidean weights above and
polynomial test functions at the stated \(N\) (for the regulated Lorentzian measure
the weight is not nonnegative, so the even-odd identity gives no sign); they do not
treat the signed fermionic Pfaffian at \(N\ge3\), which is not a nonnegative weight
and where the determinant identity is unavailable, the large-\(N\) singlet algebra
of relaxation (i) of Section~\ref{sec:status-D}, other involutions, non-polynomial
test algebras, the thimble contour integrals of~\cite{ChouNishimuraTripathi2025}, or
the slice bridge left open above. The bearing is to bound the reconstruction route
of that relaxation at the finite-\(N\) Euclidean level, complementing the
contour-side no-go, and to derive nothing: no signature, no background and no frame
is obtained, Assumption~\ref{ass:lorentzian} is not discharged. Problem~E remains
open and PROGRAMME-CONDITIONAL.
